\section{Cross-Layer Coverage Gaps}
\label{sec:gaps}

The surveyed approaches use different behavioral models and cannot be combined
by simply placing their guarantees in one list.  A cross-layer execution
argument must preserve the assumptions needed when an output expressed under
one model becomes an input to another.  Five gaps recur across the reviewed
corpus.

\subsection{Semantic and Physical Validity}

Mission contracts can make goals and contexts explicit, while planning and
learning systems can produce assignments or actions.  The receiving component
still needs the task identifiers, capabilities, map, state, dynamics, resources,
and communication model behind the proposal.  It should also be able to tell a
hard constraint from a planner preference, a model suggestion, and a checked
feasibility claim.  A meaningful plan may be physically infeasible; a feasible
plan may still violate the operator's intent.  Feasibility must therefore be
qualified as nominal, robust, probabilistic, or conditional on a named model and
operating region
\cite{mallozzi2020crome,mallozzi2022repair,mandi2023roco,obata2024lipllm,
chaki2013probabilistic}.

\subsection{Temporal Validity of Distributed State}

Distributed execution introduces several non-equivalent times: the time of the
physical event, publication time, arrival time, consumer-availability time,
forecast target time, and validity interval.  Replacing the event time with a
reception timestamp hides producer and network delay.  Age of information
measures how old the newest available update is at the consumer
\cite{kaul2012aoi}.  Transport reliability or a cadence notification does not
establish that a state is recent enough for the current decision.  A cross-robot interface therefore needs
event-time semantics, clock-domain and clock-error information, an
application-specific age or validity bound, and a response to stale or
incomparable state \cite{omg2015dds,kronauer2021latency}.

\subsection{Uncertainty, Provenance, and Evidence}

The word \emph{uncertainty} may refer to a covariance, interval, reachable set,
calibrated probability, or semantic ambiguity.  These representations cannot be
compared or propagated in the same way.  A confidence value with no declared
meaning cannot satisfy a downstream chance constraint or robust margin.  The
interface must identify the representation, covered variables and correlations,
calibration regime, validity horizon, and propagation rule.

The same requirement applies to evidence.  A consumer should be able to
distinguish a proof, solver status, monitor result, calibration report, test
identifier, and deployment measurement.  Provenance links the payload and its
evidence to the model, configuration, observation, and transformation that
produced them.  Without this lineage, downstream rejection may reveal that a
claim failed but not which assumption or artifact should be repaired
\cite{w3c2013prov}.

\subsection{Implementation and Lifecycle Compatibility}

Type and wire compatibility are weaker than behavioral compatibility.  Two
components can decode the same message while using different coordinate frames,
units, parameter sets, middleware implementations, dependency versions, or
clock sources.  Accepted execution therefore requires separate identities for
schema, protocol, implementation, dependencies, and configuration, together
with evidence for the supported combination
\cite{macenski2022ros2,omg2015dds,kronauer2021latency}.

Message contracts also omit lifecycle unless lifecycle is represented
explicitly.  Join, leave, restart, and partition recovery can invalidate old
assignments or duplicate resource ownership.  A multi-robot contract needs a
team epoch, membership and ownership state, legal lifecycle transitions, and a
reconciliation rule for retained or delayed data
\cite{luckcuck2022compositional}.

\subsection{Intervention, Recovery, and Attribution}

Monitoring is incomplete if a failed check has no defined consequence.  The
responsible authority may request clarification, obtain fresh state, re-estimate,
replan, transfer control, degrade the mission, isolate a robot, stop, or escalate
to a human.  In a team, this response must be coordinated: an unreported local
safety override can invalidate peer trajectories or task completion logic.
Runtime-assurance and field-testing work consequently treat intervention,
observability, and reporting as first-class mechanisms
\cite{desai2019soter,shivakumar2020soterros,caldas2024rvft}.

Fault origin and first observation must also be separated.  A collision risk
first detected by a controller may originate in semantic grounding, prediction,
physical modeling, communication delay, or hardware.  Recording origin,
observation, effect, and recovery authority as separate fields supports both
online handling and post-mission diagnosis.  These five gaps determine what an
edge contract must carry (\cref{sec:edge-contracts}).
